% !TeX spellcheck = it_IT
\chapter{Sviluppo e risultati sperimentali}\label{chapter:MineStuff}
Dopo aver presentato nel Capitolo precedente le nozioni 
fondamentali e le operazioni di pre-processing necessarie 
alla gestione dei dati, in questo capitolo 
verrà descritto lo sviluppo pratico e le analisi sperimentale svolta.  
L'obiettivo principale è quello di progettare e realizzare 
una metodologia computazionale in grado di 
identificare \textit{clustering alternativi} nei dati di trascrittomica spaziale, 
fornendo una rappresentazione più flessibile e interpretabile dei 
domini genici rispetto agli approcci convenzionali.

Per fare ciò è necessario che ogni gene venga e suddiviso spazialmente in sottogeni, affinchè essi si possano 
catalogare tra di loro e poter analizzare eventuali relazioni particolari che ci sono fra un sottogene e un altro.
\section{Grafo dei vicini}\label{subsec:Vicini}
Come prima cosa, come mostrato nella figura \ref{fig:five}, è necessario individuare per ogni spot $s_i$ 
l'insieme dei suoi spot adiacenti, ovvero la lista dei \textit{vicini}. 
Questa informazione è fondamentale poiché consente di stabilire quando due geni si trovano in prossimità spaziale,
condizione necessaria per poter poi definire relazioni di vicinanza, contiguità e sovrapposizione tra geni attivi.

I dati forniti dalla tecnologia 10x Visium contengono, per ciascuno spot, 
le sue coordinate sulla griglia di cattura. 
Tuttavia, tali coordinate sono organizzate in forma tabellare, come visto nella tabella \ref{tab:tsvTable}, all'interno della struttura \texttt{scdata.obs},
e perciò non immediatamente interpretabili in termini spaziali. 
Per poter costruire un grafo dei vicini è quindi necessario effettuare una conversione preliminare di queste informazioni,
trasformando le coordinate della griglia in coordinate cartesiane.
In questo modo è possibile rappresentare ogni spot come un punto nel piano, preservando le distanze e la disposizione
esagonale della griglia Visium.

\begin{lstlisting}[style=mypython,
                   label=code:index_to_cart,
                   caption={Trasformazione delle coordinate spaziali in coordinate cartesiane}]
    index_to_cart = np.zeros((len(scdata.obs), 2))
    max_row       = scdata.obs['array_row'].max()
    for i, ind in enumerate(scdata.obs.index):
        col = scdata.obs['array_col'][ind]
        row = scdata.obs['array_row'][ind]
        index_to_cart[i, 0] = col * (np.sqrt(3)/2)
        index_to_cart[i, 1] = (max_row - row) * (3/2)
\end{lstlisting}


Il frammento di codice \ref{code:index_to_cart} mostra come vengono calcolate e salvate le coordinate cartesiane di ciascuno spot,
partendo dalle colonne \texttt{array\_row} e \texttt{array\_col} della tabella di input. 
Viene inoltre applicata una correzione di scala per tenere conto della disposizione esagonale della griglia,
in cui le colonne sono sfalsate e le distanze orizzontali e verticali non sono equivalenti.

\paragraph{Costruzione del grafo dei vicini con k\text{-}NN.}
Definite le coordinate cartesiane $\mathbb{R}^2$ per ciascuno spot, 
il passo successivo consiste nell'identificare, per ogni spot $s_i$ 
i sei spot più prossimi nella griglia esagonale 10x Visium. 
A tal fine è stato impiegato l'algoritmo dei \textit{Nearest Neighbors} (k\text{-}NN), 
nella sua implementazione fornita dalla libreria \texttt{sklearn.neighbors}, 
libreria che implementa la funzione (\texttt{NearestNeighbors( ... )}).
L'idea dietro la funzione è la seguente: dato l'insieme di punti $\{x_1, ..., x_n\} \in \mathbb{R}^2$, 
per ciascun punto $\mathbf{x}_i$ si ricercano i $k$ vicini
più prossimi secondo una metrica prefissata, come la distanza euclidea. 
Poiché la funzione restituisce anche il punto stesso 
come primo vicino, si imposta $k = 1+6$ così da ottenere, 
oltre a $s_i$, i sei vicini geometrici desiderati. 

In pratica, dopo aver eseguito l'addestramento con (\texttt{NearestNeighbors( ... ).fit(coords)}), 
si invoca \texttt{kneighbors(coords)} per ottenere, per ogni riga $i$, 
le coppie \texttt{(dists[i], inds[i])}, dove \texttt{inds[i]} sono gli indici 
dei $k$ punti più vicini a $s_i$ e \texttt{dists[i]} le corrispondenti distanze.
Scartando la prima posizione, \texttt{inds[i, 0]}, che coincide con $i$, 
si ottiene la lista ordinata dei sei vicini di $s_i$.
Inoltre su ogni distanza viene fatto un confronto per verificare che sia minore o uguale del raggio prefissato, così 
da poter escludere i vicini che sono stati erroneamente ritornati dalla funzione.

Operativamente, la procedura si articola come segue: (i) si prepara la matrice \texttt{coords} di dimensione
$(n_{\text{spots}},\,2)$ contenente le coordinate cartesiane; (ii) si istanzia l’oggetto \texttt{NearestNeighbors}
specificando il numero di vicini $k$ e la metrica (\texttt{metric='euclidean'}); (iii) si addestra l’indicizzatore con
\texttt{fit}; (iv) si calcolano per tutti i punti i $k$ vicini con \texttt{kneighbors}, ottenendo indici e distanze;
(v) per ciascuno spot si selezionano i sei vicini effettivi come \texttt{inds[i, 1:]} e si memorizzano, se necessario,
anche le distanze \texttt{dists[i, 1:]} a supporto di eventuali pesature o controlli di qualità.

Dal punto di vista algoritmico, \texttt{NearestNeighbors} costruisce internamente una struttura di ricerca (KD-tree o
Ball-tree, a seconda dei dati, con selezione automatica) che consente query efficienti nello spazio bidimensionale: il
costo atteso è sub-quadratico e risulta adatto al numero di spot tipico dei vetrini Visium. L’uso della metrica euclidea
è coerente con la proiezione cartesiana effettuata in precedenza e preserva le relazioni di prossimità indotte dalla
geometria esagonale.

\paragraph{Accortezze pratiche e robustezza.}
La costruzione del grafo k\text{-}NN presenta alcune attenzioni utili in ambito sperimentale. Innanzitutto, nei bordi
della griglia reale possono mancare spot (ad esempio per ritagli o aree escluse); la scelta k-fissa garantisce comunque
sei vicini “geometrici” tra quelli effettivamente presenti, ma è consigliabile simmetrizzare il grafo finale (se $j$ è
vicino di $i$, rendere $i$ vicino di $j$) per ottenere un grafo non orientato e più stabile. In presenza di distanze
quasi uguali (tie), l’ordinamento restituito dall’algoritmo è deterministico rispetto alla struttura interna, e la
simmetrizzazione riduce sensibilmente l’effetto di scelte arbitrarie locali. Inoltre, salvare anche le distanze consente
di diagnosticare outlier geometrici (vicini insolitamente lontani) ed eventualmente applicare un filtro a raggio
(\textit{radius cutoff}) come controllo secondario.

Infine, la scelta k=6 è motivata dalla topologia esagonale di Visium: in una tessellazione ideale ciascun esagono ha sei
adiacenze. In pratiche reali con “buchi” o artefatti di imaging, una variante a raggio fisso (\texttt{radius\_neighbors})
può essere utilizzata per confermare la consistenza locale delle connessioni; tuttavia, la strategia k\text{-}NN risulta
più semplice e riproducibile per la costruzione iniziale del grafo dei vicini.

\paragraph{Estratto di codice.}
Il frammento seguente riassume l’implementazione descritta:
\begin{lstlisting}[style=mypython,caption={Ricerca dei 6 vicini più prossimi con k\text{-}NN},label=code:knn_neighbors,float,floatplacement=H]
    coords = index_to_cart  # shape: (n_spot, 2)
    k = 1 + 6 
    nbrs = NearestNeighbors(n_neighbors=k, metric='euclidean').fit(coords)
    dists, inds = nbrs.kneighbors(coords)

    i = 0  # ad esempio: i = indice dello spot
    vicini6  = inds[i, 1:]
    distanze = dists[i, 1:]
\end{lstlisting}

A valle di questo passaggio, la struttura \texttt{inds} costituisce la base per costruire l’elenco di adiacenze del
grafo dei vicini, che verrà poi utilizzato nelle analisi di connettività (componenti connesse), nelle misure di
vicinanza tra geni attivi e, più in generale, in tutte le operazioni che richiedono consapevolezza della geometria
spaziale del tessuto.
